\documentclass[12pt]{article}
\usepackage[T1]{fontenc}
\usepackage[utf8]{inputenc}
\usepackage[brazil]{babel}
\usepackage[margin=1in]{geometry}
\usepackage{ragged2e}
\usepackage[normalem]{ulem}
\setlength{\parindent}{0pt}
\setlength{\parskip}{0.8\baselineskip}
\begin{document}
\justifying
\noindent Verifica-se que a questão objeto da controvérsia jurídica suscitada no recurso manejado pela parte recorrente esteve afetada no representativo da controvérsia - \textbf{Tema n. }\textbf{1} da TNU - (PU no processo n. PEDILEF 2009.51.51.066212-3/ RJ),  afetado em 10/06/2011 foi julgado (trânsito em julgado em 05/10/2011), por meio do qual foi elucidada a seguinte questão:

\noindent \textit{"Saber qual a forma de cálculo da aposentadoria por invalidez, do auxílio-doença e pensões derivadas, em momento antecedente à edição da Lei n. 9.876/99."}

\noindent Restou firmada a seguinte tese:

\noindent \textit{"O valor da aposentadoria por invalidez e do auxílio-doença, bem como das pensões destes derivados ou calculadas com base no art. 75, da Lei n. 8.213/91, será obtido, na forma do art. 29, II, do mesmo diploma, por meio da média aritmética simples dos 80\% maiores salários de contribuição, considerado todo o período contributivo, independentemente do momento de inscrição do segurado e do número de contribuições mensais do período contributivo." (Tema 1 TNU).}

\noindent Com efeito, depreende-se da leitura do voto condutor do acórdão que o mesmo se encontra em conformidade com o decidido no  tema pela TNU.

\noindent A proposito, consta do acordao hostilizado: (...) Indenização por dano moral, Responsabilidade civil, DIREITO CIVIL

\noindent Nessas condições, o conteúdo do Acórdão recorrido é consentâneo com o entendimento da TNU, firmado em julgamento de recurso representativo de controvérsia, o que atrai a incidência da regra prevista pelo art. 14, III, \textit{b,} da Resolução CJF n. 586/2019:

\noindent \textit{\textbf{Art. 14.}}\textit{ Decorrido o prazo para contrarrazões, os autos serão conclusos ao magistrado responsável pelo exame preliminar de admissibilidade, que deverá, de forma sucessiva: (...) }\textit{\textbf{III }}\textit{- negar seguimento a pedido de uniformização de interpretação de lei federal interposto contra acórdão que esteja em conformidade com entendimento consolidado: }\textit{\textbf{b)}}\textit{ em }\uline{\textit{\textbf{recurso representativo de controvérsia pela Turma Nacional de Uniformização}}}\textit{ ou em pedido de uniformização de interpretação de lei dirigido ao Superior Tribunal de Justiça;(...).}

\noindent Posto isso, com arrimo no \uline{\textbf{art. 14, III, }}\uline{\textit{\textbf{b}}}\textit{,} da Resolução CJF n. 586/2019 c/c art. 1.030, I, do CPC, \textbf{NEGO SEGUIMENTO }\textbf{a}o\textbf{ PU}\textbf{.}

\noindent Intimem-se. Aguarde-se o decurso do prazo recursal. Após, certifique-se o trânsito em julgado e devolva-se ao juízo de origem.
\end{document}
